% -*- coding: utf-8 -*-
% Documentación Técnica: Modelos Fundacionales e Infraestructura de Servidor GPU
% Trabajo de Fin de Grado — FatigueSet
% Autor: Emilio Gullón
% Fecha: Julio 2026

\documentclass[12pt, a4paper]{article}

% ---- Paquetes ----
\usepackage[utf8]{inputenc}
\usepackage[T1]{fontenc}
\usepackage[spanish]{babel}
\usepackage{amsmath, amssymb}
\usepackage{graphicx}
\usepackage{booktabs}
\usepackage{longtable}
\usepackage{listings}
\usepackage{xcolor}
\usepackage{hyperref}
\usepackage{geometry}
\usepackage{caption}
\usepackage{subcaption}
\usepackage{float}
\usepackage{enumitem}
\usepackage{csquotes}
\usepackage{multirow}

\geometry{margin=2.5cm}
\hypersetup{
    colorlinks=true,
    linkcolor=blue!70!black,
    citecolor=green!50!black,
    urlcolor=blue!60!black,
}

% ---- Estilo de código Python ----
\definecolor{codebg}{HTML}{F7F7F7}
\definecolor{codegreen}{HTML}{28A745}
\definecolor{codeorange}{HTML}{D73A49}
\definecolor{codeblue}{HTML}{0366D6}
\definecolor{codegray}{HTML}{6A737D}

\lstdefinestyle{pythonstyle}{
    language=Python,
    backgroundcolor=\color{codebg},
    basicstyle=\ttfamily\small,
    keywordstyle=\color{codeorange}\bfseries,
    stringstyle=\color{codegreen},
    commentstyle=\color{codegray}\itshape,
    identifierstyle=\color{codeblue},
    numbers=left,
    numberstyle=\tiny\color{codegray},
    numbersep=8pt,
    frame=single,
    rulecolor=\color{codegray!30},
    breaklines=true,
    showstringspaces=false,
    tabsize=4,
    captionpos=b,
    xleftmargin=1.5em,
}
\lstset{style=pythonstyle}

% ---- Comandos útiles ----
\newcommand{\code}[1]{\texttt{#1}}
\newcommand{\fatigueset}{\textsc{FatigueSet}}
\newcommand{\moment}{\textsc{MOMENT}}
\newcommand{\chronos}{\textsc{Chronos}}
\newcommand{\timesfm}{\textsc{TimesFM}}

% ---- Metadatos ----
\title{
    \vspace{-2cm}
    \textbf{Documentación Técnica: Modelos Fundacionales\\
    para Series Temporales e Infraestructura GPU} \\[0.5em]
    \large Trabajo de Fin de Grado --- Predicción de Fatiga con \fatigueset{}
}
\author{Emilio Gullón}
\date{Julio 2026}

\begin{document}
\maketitle

\begin{abstract}
Este documento recoge la implementación, despliegue y evaluación experimental de tres modelos fundacionales para series temporales --- \moment{} (encoder-only), \chronos{} (decoder-only probabilístico) y \timesfm{} (decoder-only con parches) --- aplicados a la predicción de fatiga física y mental a partir de señales fisiológicas del dataset \fatigueset{}. Se describe en detalle la librería \code{fatigueset.models.foundation}, el módulo de optimización de hiperparámetros con Optuna (\code{fatigueset.models.optimizers}), la infraestructura de ejecución en un clúster GPU universitario (UGR), y los resultados comparativos frente a modelos clásicos de Machine Learning y redes de Deep Learning diseñadas a medida. Adicionalmente, se documentan los problemas técnicos encontrados (errores de memoria GPU, límites de tiempo en SLURM) y las soluciones aplicadas.
\end{abstract}

\tableofcontents
\newpage

% ======================================================================
\section{Introducción}
\label{sec:introduccion}
% ======================================================================

\subsection{Contexto del TFG y FatigueSet}

Este Trabajo de Fin de Grado aborda la predicción cuantitativa de fatiga física y mental a partir de señales fisiológicas multicanal del dataset \fatigueset{} \cite{kalanadhabhatta2022fatigueset}. El dataset contiene registros de 12 participantes monitorizados con dispositivos portátiles (Empatica E4 y Zephyr BioHarness) durante tareas cognitivas y físicas, proporcionando 23 canales fisiológicos (acelerómetros, electrodermal, temperatura, ritmo cardíaco, respiración, etc.) junto con puntuaciones subjetivas de fatiga en escala 0--100.

El pipeline experimental se estructura en cuatro fases progresivas de complejidad:
\begin{enumerate}[label=\arabic*.]
    \item \textbf{Modelos clásicos de Machine Learning}: Random Forest, KNN, Lasso, Ridge, SVM, Decision Tree, ElasticNet y Regresión Lineal.
    \item \textbf{Modelos de Deep Learning diseñados a medida}: LSTM, GRU, CNN-LSTM, TCN, Transformer, PatchTST y xLSTM.
    \item \textbf{Optimización de hiperparámetros con Optuna}: Búsqueda bayesiana con TPE (Tree-structured Parzen Estimator) y optimizadores personalizados.
    \item \textbf{Modelos fundacionales preentrenados}: \moment{}, \chronos{} y \timesfm{}.
\end{enumerate}

\subsection{Motivación: modelos fundacionales para series temporales}

Los modelos fundacionales (\textit{Foundation Models}) representan un cambio de paradigma en el análisis de series temporales. Frente al enfoque tradicional de diseñar y entrenar un modelo específico para cada tarea, estos modelos se preentrenan sobre grandes corpora de series temporales diversas y posteriormente se adaptan a tareas downstream con mínima supervisión.

El survey de Kottapalli et al. \cite{kottapalli2025survey} identifica tres paradigmas de adaptación:
\begin{itemize}
    \item \textbf{Zero-shot}: el modelo preentrenado se aplica directamente sin ningún reentrenamiento.
    \item \textbf{Fine-tuning}: se reentrenan las últimas capas o una cabeza de tarea ligera.
    \item \textbf{Prompting / In-context learning}: se proporcionan ejemplos de demostración al modelo.
\end{itemize}

En este TFG se aplican los tres modelos fundacionales más relevantes identificados en el survey, evaluando su capacidad como extractores de características universales para la predicción de fatiga en un dominio no visto durante su preentrenamiento.

\subsection{Objetivos de esta fase experimental}

Los objetivos específicos de la fase de modelos fundacionales son:
\begin{enumerate}[label=(\alph*)]
    \item Implementar adaptadores para \moment{}, \chronos{} y \timesfm{} dentro de la librería modular \code{fatigueset}.
    \item Evaluar el rendimiento de cada modelo en predicción de fatiga física y mental usando validación cruzada GroupKFold por participante.
    \item Comparar métricas deterministas (MAE, RMSE, $R^2$) y probabilísticas (CRPS, cobertura del intervalo de predicción al 90\%).
    \item Desplegar y ejecutar los experimentos a escala completa en un clúster GPU universitario.
\end{enumerate}


% ======================================================================
\section{Marco Teórico: Modelos Fundacionales para Series Temporales}
\label{sec:marco_teorico}
% ======================================================================

\subsection{Taxonomía de arquitecturas}

Siguiendo la taxonomía del survey de Kottapalli et al. \cite{kottapalli2025survey}, los modelos fundacionales para series temporales se clasifican según su arquitectura base:

\begin{itemize}
    \item \textbf{Encoder-only}: procesan la secuencia completa bidireccionalmente para producir representaciones contextualizadas. Ejemplo: \moment{}.
    \item \textbf{Decoder-only}: procesan la secuencia de forma autorregresiva (causal), prediciendo el siguiente token/valor. Ejemplos: \chronos{}, \timesfm{}.
    \item \textbf{Encoder-Decoder}: combinan ambos enfoques. Ejemplo: Lag-Llama (no implementado en este TFG).
\end{itemize}

La Tabla~\ref{tab:taxonomia_modelos} resume la taxonomía de los modelos implementados.

\begin{table}[H]
\centering
\caption{Taxonomía de los modelos fundacionales implementados en el TFG.}
\label{tab:taxonomia_modelos}
\small
\begin{tabular}{@{}lccccc@{}}
\toprule
\textbf{Modelo} & \textbf{Tipo} & \textbf{Patch-based} & \textbf{Multivariado} & \textbf{Probabilístico} & \textbf{Paradigma} \\
\midrule
\moment{} & Encoder-only & Sí & Sí & No & Fine-tuning cabeza \\
\chronos{} & Decoder-only & No & No & Sí & Zero-shot + probe \\
\timesfm{} & Decoder-only & Sí & No & Sí & Zero-shot + probe \\
\bottomrule
\end{tabular}
\end{table}


\subsection{MOMENT: arquitectura, parches y preentrenamiento}

\moment{} (\textit{A Family of Open Time-series Foundation Models}) \cite{goswami2024moment} es un modelo encoder-only basado en Transformer preentrenado mediante reconstrucción de parches enmascarados (\textit{Masked Patch Reconstruction}), análogo al preentrenamiento de BERT en NLP pero adaptado a series temporales.

\paragraph{Arquitectura.}
La secuencia temporal de entrada $\mathbf{x} \in \mathbb{R}^{T \times C}$ (con $T$ pasos temporales y $C$ canales) se segmenta en \textit{parches} no superpuestos de longitud fija $P$. Cada parche se proyecta a un espacio de dimensión $d_{\text{model}}$ mediante una capa lineal, produciendo una secuencia de tokens:
\[
\mathbf{z}_i = W_{\text{patch}} \cdot \mathbf{x}_{[iP:(i+1)P]} + \mathbf{b}_{\text{patch}}, \quad i = 0, 1, \ldots, \lfloor T/P \rfloor - 1
\]

Los tokens se procesan por $L$ capas de Transformer encoder (auto-atención multi-cabeza + MLP feed-forward), produciendo representaciones contextualizadas $\mathbf{h}_i$ para cada parche.

\paragraph{Preentrenamiento.}
\moment{} se preentrena en el corpus \textit{Time Series Pile}, una colección de más de 1 billón de observaciones de series temporales de dominios diversos (finanzas, salud, clima, energía). El objetivo de preentrenamiento es la reconstrucción de parches enmascarados: se enmascara un porcentaje aleatorio de parches y el modelo debe reconstruirlos a partir del contexto.

\paragraph{Variantes.} Se implementan dos variantes:
\begin{itemize}
    \item \code{MOMENT-1-small}: 40M parámetros (pruebas locales).
    \item \code{MOMENT-1-large}: 385M parámetros (experimento en servidor, 341.250.570 parámetros totales).
\end{itemize}

\paragraph{Adaptación a FatigueSet.}
Para adaptar \moment{} a la predicción de fatiga, se congela el backbone preentrenado y se añade una cabeza de regresión lineal \code{Linear($d_{\text{model}}$, 2)} que mapea la representación agregada (media de los embeddings de todos los parches) a las dos dimensiones de fatiga. Solo se entrenan los 2.050 parámetros de la cabeza (0.00\% del total), lo que reduce drásticamente el riesgo de sobreajuste en el pequeño dataset de \fatigueset{}.


\subsection{Chronos: tokenización por quantile binning y T5}

\chronos{} \cite{ansari2024chronos} es un framework probabilístico decoder-only para series temporales que reutiliza la arquitectura T5 (\textit{Text-to-Text Transfer Transformer}) de Google, adaptándola al dominio temporal mediante una estrategia de tokenización numérica.

\paragraph{Tokenización.}
Los valores continuos de la serie temporal se discretizan en \textit{bins} mediante \textit{quantile binning}: se calculan los cuantiles empíricos del contexto y cada valor se asigna al bin correspondiente. Esta estrategia permite reutilizar el vocabulario discreto del T5 sin modificar la arquitectura.

\paragraph{Preentrenamiento.}
\chronos{} se preentrena con un objetivo de lenguaje estándar (predicción del siguiente token) sobre el corpus TSMixup, una mezcla de series temporales reales y sintéticas generadas por Gaussian Processes.

\paragraph{Variantes.}
\begin{itemize}
    \item \code{chronos-t5-tiny}: 8M parámetros (pruebas locales).
    \item \code{chronos-t5-large}: 710M parámetros (experimento en servidor).
\end{itemize}

\paragraph{Adaptación a FatigueSet.}
Dado que \chronos{} está preentrenado exclusivamente para \textit{forecasting} univariado, la adaptación a regresión multidimensional de fatiga se realiza mediante un enfoque de \textbf{linear probe} (sonda lineal):
\begin{enumerate}
    \item Cada uno de los 23 canales fisiológicos se trata como una serie temporal univariada independiente.
    \item \chronos{} predice la distribución del siguiente paso para cada canal, obteniéndose la mediana como característica.
    \item Las 23 medianas se concatenan como vector de características y se alimentan a una regresión Ridge que mapea al espacio de fatiga bidimensional.
\end{enumerate}


\subsection{TimesFM: parches de longitud variable y decoder patched}

\timesfm{} 2.5 (\textit{Time Series Foundation Model}) \cite{das2024timesfm} es un modelo decoder-only de Google Research basado en parches de longitud variable. A diferencia de \moment{}, que usa parches de longitud fija, \timesfm{} utiliza un mecanismo de \textit{input patching} con múltiples resoluciones que permite procesar series temporales de diferentes frecuencias sin re-muestreo.

\paragraph{Arquitectura.}
\timesfm{} 2.5 utiliza un decoder Transformer con 200M de parámetros, entrenado sobre 100 billones de observaciones de series temporales reales y sintéticas. Su diseño clave es la predicción por parches con tamaño de salida variable (\textit{output patch size}), lo que lo hace extremadamente eficiente en inferencia.

\paragraph{Adaptación a FatigueSet.}
La adaptación sigue el mismo paradigma de channel independence y linear probe que \chronos{}:
\begin{enumerate}
    \item Los 23 canales se procesan independientemente por \timesfm{}.
    \item El modelo genera predicciones puntuales y distribuciones de cuantiles (10 cuantiles predefinidos) por canal.
    \item Las predicciones puntuales se concatenan como vector de 23 características y se alimentan a una regresión Ridge.
    \item Los cuantiles se usan para calcular métricas probabilísticas (CRPS, cobertura).
\end{enumerate}


% ======================================================================
\section{Modelos Clásicos de Aprendizaje Automático}
\label{sec:modelos_clasicos}
% ======================================================================

Esta sección documenta los experimentos realizados con algoritmos de aprendizaje automático clásico (Machine Learning clásico). Estos modelos se entrenaron sobre un dataset estructurado tabular obtenido mediante agregación temporal de las señales fisiológicas.

\subsection{Protocolo Experimental y Preprocesamiento}

Para los modelos clásicos, se estructuró el dataset tabular \code{df\_ml} que consta de 108 filas (correspondientes a los 12 participantes, evaluados en 3 sesiones independientes con 3 fases experimentales cada una). Cada muestra contiene más de 50 características agregadas que resumen los estadísticos (media, desviación estándar) de los sensores cardíacos (frecuencia cardíaca, variabilidad cardíaca, respiración), sensores periféricos de la muñeca (temperatura, conductancia de la piel) y bandas de potencia de EEG.

La evaluación se realizó utilizando un protocolo de validación cruzada de 5 particiones (\code{KFold} aleatorio, $k=5$, con semilla fija para reproducibilidad). Se entrenaron los siguientes regresores de la librería \code{scikit-learn}:
\begin{itemize}
    \item \textbf{Random Forest Regressor}: Ensamble por \textit{bagging} con 50 árboles de decisión y profundidad máxima de 10.
    \item \textbf{K-Nearest Neighbors (KNN)}: Regresor basado en vecindad con pesos uniformes ($k=3$ y $k=5$).
    \item \textbf{Lasso Regressor}: Regresión lineal con regularización $L_1$ ($\alpha=0.1$).
    \item \textbf{Decision Tree Regressor}: Árbol de decisión simple sin límite de profundidad (utilizado como baseline de sobreajuste).
    \item \textbf{ElasticNet}: Regresión lineal regularizada con penalización combinada $L_1$ y $L_2$ ($\alpha=0.1$).
    \item \textbf{Ridge Regressor}: Regresión lineal regularizada con penalización $L_2$ ($\alpha=1.0$).
    \item \textbf{Support Vector Regressor (SVR)}: Máquina de vectores de soporte con kernel de función de base radial (RBF) y constante de penalización $C=100$.
    \item \textbf{Regresión Lineal}: Regresión por mínimos cuadrados ordinarios sin regularización.
\end{itemize}

\subsection{Ranking y Comparativa de Rendimiento}

Los modelos se evaluaron usando $R^2$ de validación cruzada y el MAE medio sobre el conjunto completo de entrenamiento. En la Tabla~\ref{tab:ranking_clasicos} se muestra el ranking de los modelos para la fatiga física.

\begin{table}[H]
\centering
\caption{Ranking y comparación de modelos clásicos para Fatiga Física (CV de 5 splits).}
\label{tab:ranking_clasicos}
\small
\begin{tabular}{@{}lccccc@{}}
\toprule
\textbf{Modelo} & \textbf{Tipo} & \textbf{$R^2$ CV Medio} & \textbf{Desv. Est. $R^2$ CV} & \textbf{MSE Train} & \textbf{MAE Train} \\
\midrule
Random Forest & Ensamble & $-$2.4084 & 2.6050 & 1.0772 & 0.8767 \\
KNN (k=3) & Instancias & $-$3.5707 & 3.2852 & 4.8593 & 1.6889 \\
Lasso ($\alpha=0.1$) & Regularizado & $-$3.8720 & 3.7321 & 8.2614 & 2.3716 \\
Decision Tree & Árbol & $-$4.7405 & 6.6825 & 0.0000 & 0.0000 \\
KNN (k=5) & Instancias & $-$4.8633 & 3.4261 & 8.5387 & 2.3467 \\
ElasticNet & Regularizado & $-$5.7379 & 7.3465 & 11.3750 & 2.9468 \\
Ridge ($\alpha=1.0$) & Regularizado & $-$6.0179 & 7.9491 & 11.3456 & 2.9367 \\
SVR & Kernel & $-$6.5572 & 8.6872 & 7.6892 & 2.1611 \\
Regresión Lineal & Lineal & $-$12.5545 & 15.5509 & 5.5752 & 1.9046 \\
\bottomrule
\end{tabular}
\end{table}

\subsection{Análisis del Impacto de la Normalización por Sujeto (Subject-wise Normalization)}
\label{subsec:normalizacion_sujeto}

Una de las propuestas de preprocesamiento analizadas para mitigar la heterogeneidad individual (diferencias de línea base y rango dinámico de las señales fisiológicas entre participantes) fue la aplicación de una \textbf{Normalización por Participante y Sesión (Subject-Session-wise Normalization)}.

\paragraph{Formulación Matemática.}
Para cada participante $i$ y sesión $s$, cada característica $x$ se escala de forma independiente mediante Z-score utilizando únicamente la media $\mu_{i,s}$ y la desviación estándar $\sigma_{i,s}$ de ese bloque de datos:
\[
z_{i,s,t} = \frac{x_{i,s,t} - \mu_{i,s}}{\sigma_{i,s}}
\]
Donde $t$ representa los pasos temporales dentro de la sesión.

\paragraph{Evaluación Empírica.}
Para cuantificar rigurosamente su impacto, se realizó un experimento comparativo utilizando un modelo robusto de Random Forest bajo validación cruzada GroupKFold por participante (garantizando que ningún dato de los participantes de test se use durante el entrenamiento). Los resultados obtenidos para la predicción de fatiga física son:
\begin{itemize}
    \item \textbf{Sin Normalización}: $\text{MAE} = 2.6878$ | $R^2 = 0.4120$
    \item \textbf{Con Normalización por Sujeto/Sesión}: $\text{MAE} = 3.9222$ | $R^2 = -0.2614$
\end{itemize}

El uso de la normalización por participante y sesión provocó un incremento del \textbf{45.9\% en el error absoluto medio (MAE)} y redujo el coeficiente de determinación $R^2$ de una puntuación muy positiva (0.41) a una negativa ($-0.26$), destruyendo la capacidad predictiva del modelo.

\paragraph{Justificación del Empeoramiento.}
El fracaso de esta normalización se debe a tres factores clave:
\begin{enumerate}
    \item \textbf{Pérdida de la varianza inter-sesión (línea base de fatiga)}: Al centrar a cero cada sesión por separado, se elimina la diferencia de escala absoluta entre una sesión en la que el sujeto está relajado (frecuencia cardíaca baja) y una sesión de esfuerzo o fatiga (frecuencia cardíaca alta). Si ambas sesiones se escalan para tener media cero, sus representaciones se vuelven casi idénticas a ojos del modelo, imposibilitando predecir la diferencia de fatiga.
    \item \textbf{Amplificación del ruido}: En sesiones de muy baja actividad o estados muy estables, la desviación estándar $\sigma_{i,s}$ es extremadamente pequeña. Al dividir por un valor cercano a cero, cualquier ruido térmico o pequeña fluctuación de lectura del sensor se magnifica en el Z-score, actuando como ruido severo.
    \item \textbf{Pérdida de leyes biológicas universales}: Las respuestas físicas humanas ante la fatiga siguen tendencias absolutas (p. ej., a mayor fatiga, mayor ritmo cardíaco y mayor conductancia galvánica de la piel). La normalización local descarta esta escala absoluta común, dificultando que el modelo generalice sus conocimientos a nuevos sujetos invisibles.
\end{enumerate}

\subsection{Importancia de Características (Feature Importance)}

El modelo \textbf{Random Forest Regressor} permite extraer la importancia relativa de las variables basándose en la reducción de impureza (Gini/varianza). La Tabla~\ref{tab:importancia_features} presenta las 15 características fisiológicas más importantes identificadas por el modelo para predecir la fatiga física.

\begin{table}[H]
\centering
\caption{Top 15 características fisiológicas más importantes (Random Forest).}
\label{tab:importancia_features}
\small
\begin{tabular}{@{}llc@{}}
\toprule
\textbf{Posición} & \textbf{Característica} & \textbf{Importancia Relativa} \\
\midrule
1 & \code{eda\_media} (wrist) & 0.1319 \\
2 & \code{hr\_media} (chest) & 0.1217 \\
3 & \code{eeg\_beta\_media} (EEG) & 0.0971 \\
4 & \code{br\_media} (chest) & 0.0839 \\
5 & \code{eda\_std} (wrist) & 0.0750 \\
6 & \code{hr\_std} (chest) & 0.0679 \\
7 & \code{hrv\_media} (chest) & 0.0667 \\
8 & \code{delta\_fisica\_cognitivo} (diferencia) & 0.0566 \\
9 & \code{hrv\_std} (chest) & 0.0499 \\
10 & \code{br\_std} (chest) & 0.0442 \\
11 & \code{eeg\_alpha\_media} (EEG) & 0.0438 \\
12 & \code{eeg\_theta\_media} (EEG) & 0.0413 \\
13 & \code{delta\_mental\_cognitivo} (diferencia) & 0.0300 \\
14 & \code{delta\_mental\_ejercicio} (diferencia) & 0.0237 \\
15 & \code{mental\_M3} (estadístico) & 0.0233 \\
\bottomrule
\end{tabular}
\end{table}

\paragraph{Análisis de características.}
La alta relevancia de \code{eda\_media} y \code{hr\_media} concuerda con la teoría fisiológica de la fatiga, donde la conductancia galvánica y la frecuencia cardíaca actúan como marcadores directos del estrés y la carga cardiovascular. Asimismo, la presencia de las bandas beta y alfa de EEG evidencia la implicación del cansancio neurológico en las fases cognitivas del experimento.


% ======================================================================
\section{Modelos de Deep Learning Diseñados a Medida}
\label{sec:modelos_deep_learning}
% ======================================================================

En esta fase del proyecto, para modelar la dinámica temporal compleja de las señales fisiológicas multicanal, se diseñaron e implementaron en PyTorch siete arquitecturas de Deep Learning a medida. Estas redes se entrenaron utilizando secuencias multivariadas de timesteps correspondientes a ventanas temporales del dataset, mapeando la entrada de forma bidimensional hacia las etiquetas continuas de fatiga física y mental.

\subsection{Resumen de Arquitecturas Diseñadas}

Las arquitecturas propuestas abarcan un amplio espectro de paradigmas de procesamiento temporal:
\begin{itemize}
    \item \textbf{Custom LSTM y GRU}: Modelos recurrentes estándar para capturar dependencias secuenciales a largo plazo.
    \item \textbf{Custom CNN-LSTM}: Modelo híbrido que combina capas convolucionales 1D para la extracción de características espaciales locales y capas recurrentes LSTM para capturar la evolución temporal.
    \item \textbf{Custom TCN (Temporal Convolutional Networks)}: Arquitectura convolucional que emplea convoluciones causales y dilatadas junto con conexiones residuales, ofreciendo un campo receptivo amplio y evitando el desvanecimiento del gradiente sin la necesidad de recurrencia.
    \item \textbf{Custom Transformer}: Modelo basado en auto-atención multi-cabeza y codificación posicional para capturar relaciones globales sin restricciones secuenciales.
    \item \textbf{Custom PatchTST}: Adaptación de la arquitectura de Transformer que opera sobre parches o segmentos locales en lugar de timesteps individuales, e implementa independencia de canales (\emph{channel independence}).
    \item \textbf{Custom xLSTM (sLSTM)}: Red recurrente avanzada basada en el reciente paradigma xLSTM, implementada de forma manual para sortear las limitaciones de las compuertas sigmoideas tradicionales.
\end{itemize}

\subsection{Extended Long Short-Term Memory (xLSTM / sLSTM)}
\label{subsec:xlstm_theory}

El modelo \textbf{xLSTM} (Extended Long Short-Term Memory) \cite{beck2024xlstm} introduce mejoras fundamentales sobre la LSTM clásica mediante el uso de compuertas exponenciales y mecanismos de estabilización numérica. En este trabajo, ante la ausencia de una implementación oficial compatible en la librería estándar de PyTorch, se desarrolló manualmente una variante \textbf{sLSTM} (Stabilized LSTM) en el módulo \code{xlstm.py}.

\subsubsection{Formulación Matemática de la Celda sLSTM}

A diferencia de la LSTM clásica, que utiliza funciones sigmoideas para las compuertas de entrada y olvido (limitando su rango a $(0, 1)$), la variante sLSTM emplea funciones exponenciales para permitir que la red actualice sus estados con mayor flexibilidad. Sin embargo, el uso directo de exponenciales conlleva riesgos de desbordamiento numérico (\emph{overflow}). Para resolver esto, sLSTM introduce un estado estabilizador $m_t$ y un normalizador $n_t$.

Las ecuaciones que describen el paso hacia adelante (\emph{forward pass}) para un paso temporal $t$ en la celda sLSTM son:

\begin{enumerate}
    \item \textbf{Pre-activaciones lineales}:
    \begin{align}
        \tilde{f}_t &= x_t W_f^T + h_{t-1} U_f^T + b_f \\
        \tilde{i}_t &= x_t W_i^T + h_{t-1} U_i^T + b_i \\
        \tilde{c}_t &= \tanh(x_t W_c^T + h_{t-1} U_c^T + b_c) \\
        \tilde{o}_t &= x_t W_o^T + h_{t-1} U_o^T + b_o
    \end{align}
    Donde $\tilde{f}_t$, $\tilde{i}_t$, $\tilde{c}_t$ y $\tilde{o}_t$ son las pre-activaciones de la compuerta de olvido, compuerta de entrada, candidato de celda y compuerta de salida, respectivamente. Los pesos $W_*$ y $U_*$ y sesgos $b_*$ son parámetros aprendidos.

    \item \textbf{Estabilización en escala logarítmica}:
    Para evitar el desbordamiento numérico al calcular los exponenciales de las compuertas, se calcula un término de escala $m_t$ que mantiene el exponente máximo observado:
    \begin{equation}
        m_t = \max(m_{t-1} + \tilde{f}_t, \tilde{i}_t)
    \end{equation}
    Con una condición de inicialización $m_0 = 0$.

    \item \textbf{Compuertas Exponenciales Estabilizadas}:
    Las compuertas de olvido y de entrada se calculan en base a la escala de estabilización $m_t$:
    \begin{align}
        f_t^s &= \exp(m_{t-1} + \tilde{f}_t - m_t) \\
        i_t^s &= \exp(\tilde{i}_t - m_t)
    \end{align}
    De esta forma, la suma de los exponentes está acotada superiormente por $0$, evitando de forma robusta desbordamientos numéricos por arriba.

    \item \textbf{Actualización del Estado de Celda y Normalizador}:
    El estado interno de la celda $c_t$ y el normalizador acumulativo $n_t$ se actualizan secuencialmente:
    \begin{align}
        c_t &= f_t^s \odot c_{t-1} + i_t^s \odot \tilde{c}_t \\
        n_t &= f_t^s \odot n_{t-1} + i_t^s
    \end{align}
    Donde $\odot$ denota el producto de Hadamard (elemento a elemento). El normalizador se inicializa como $n_0 = 1$.

    \item \textbf{Normalización del Estado de Celda y Salida}:
    Para contrarrestar el crecimiento debido a los términos exponenciales, el estado de la celda se normaliza antes de proyectarlo al estado oculto:
    \begin{equation}
        \hat{c}_t = \frac{c_t}{n_t + \epsilon}
    \end{equation}
    Donde $\epsilon = 10^{-8}$ es una constante para evitar divisiones por cero. Finalmente, el estado oculto $h_t$ se obtiene como:
    \begin{align}
        o_t &= \sigma(\tilde{o}_t) \\
        h_t &= o_t \odot \tanh(\hat{c}_t)
    \end{align}
\end{enumerate}

\subsubsection{Implementación en la Biblioteca}

La celda anterior se programó en la clase \code{CustomxLSTMCell} dentro de \code{fatigueset/models/xlstm.py}. Posteriormente, la clase \code{CustomxLSTM} define la pila de capas recurrentes, permitiendo conectar múltiples capas de celdas sLSTM y aplicando regularización por \emph{dropout} entre capas intermedias. La clase final \code{CustomxLSTMRegressor} hereda de \code{nn.Module} y mapea el último paso temporal de la secuencia mediante una proyección lineal al espacio continuo de salida bidimensional de la fatiga:

\begin{lstlisting}[caption={Estructura de la clase CustomxLSTMRegressor.}, label={lst:xlstm_regressor}]
class CustomxLSTMRegressor(nn.Module):
    def __init__(self, input_size: int, hidden_size: int = 64, num_layers: int = 2, dropout: float = 0.2, output_size: int = 2):
        super().__init__()
        self.xlstm = CustomxLSTM(input_size, hidden_size, num_layers, dropout)
        self.fc = nn.Linear(hidden_size, output_size)

    def forward(self, x: torch.Tensor) -> torch.Tensor:
        # x: (batch_size, seq_len, input_size)
        out, _ = self.xlstm(x)
        last_step_out = out[:, -1, :]
        return self.fc(last_step_out)
\end{lstlisting}


% ======================================================================
\section{Implementación de la Librería \code{fatigueset.models.foundation}}
\label{sec:implementacion_foundation}
% ======================================================================

El módulo \code{foundation.py} (924 líneas, 34.6 KB) implementa toda la lógica de adaptación de los tres modelos fundacionales. Se estructura en cinco componentes principales.

\subsection{\code{MOMENTFatigueRegressor}: diseño y forward pass}

La clase \code{MOMENTFatigueRegressor} hereda de \code{torch.nn.Module} y encapsula el backbone \moment{} con una cabeza de regresión personalizada.

\begin{lstlisting}[caption={Paso forward de MOMENTFatigueRegressor (simplificado).}, label={lst:moment_forward}]
class MOMENTFatigueRegressor(nn.Module):
    def forward(self, x: torch.Tensor) -> torch.Tensor:
        # Transponer: (B, seq_len, n_channels) -> (B, n_channels, seq_len)
        x_moment = x.transpose(1, 2).contiguous()

        # Embeddings preentrenados (backbone congelado)
        outputs = self._backbone.embed(x_enc=x_moment, reduction="mean")
        pooled = outputs.embeddings  # (B, d_model)

        # Cabeza de regresion entrenada
        return self.regression_head(self.dropout_layer(pooled))
\end{lstlisting}

\paragraph{Diseño de la cabeza.} La cabeza de regresión consiste en:
\begin{itemize}
    \item \code{Dropout(p=0.1)}: regularización para evitar sobreajuste.
    \item \code{Linear($d_{\text{model}}$, 2)}: proyección lineal al espacio de fatiga bidimensional, con inicialización Xavier uniforme para los pesos y ceros para los sesgos.
\end{itemize}

\paragraph{Transposición de ejes.} \moment{} espera la entrada en formato \code{(batch, n\_channels, seq\_len)}, mientras que el resto de modelos de \fatigueset{} usan \code{(batch, seq\_len, n\_channels)}. El método \code{forward} realiza la transposición internamente, garantizando la compatibilidad sin modificar el pipeline de datos.

\paragraph{Carga perezosa del backbone.} El backbone se carga de forma perezosa (\textit{lazy loading}): si se proporciona un backbone pre-cargado en el constructor (parámetro \code{backbone}), se reutiliza directamente; en caso contrario, se carga automáticamente en el primer \code{forward} pass mediante \code{MOMENTPipeline.from\_pretrained}.


\subsection{\code{ChronosZeroShotEvaluator}: sonda lineal sobre features zero-shot}

La clase \code{ChronosZeroShotEvaluator} implementa la evaluación zero-shot de \chronos{} con tres métodos principales:

\begin{enumerate}
    \item \code{extract\_features(X)}: itera sobre los 23 canales, llama a \code{predict\_quantiles} de \chronos{} para cada canal, y extrae la mediana (cuantil 0.5) del siguiente paso como característica. Resultado: matriz \code{(N, 23)}.
    \item \code{fit\_probe(X\_train, y\_train, X\_val)}: entrena una regresión Ridge ($\alpha=1.0$) sobre las 23 características extraídas y predice sobre el conjunto de validación.
    \item \code{predict\_quantiles\_for\_crps(X)}: predice 9 cuantiles (0.1 a 0.9) por canal para el cálculo de métricas probabilísticas.
\end{enumerate}


\subsection{\code{TimesFMZeroShotEvaluator}: cuantiles y predicción por canal}

La clase \code{TimesFMZeroShotEvaluator} sigue el mismo paradigma pero aprovecha la API nativa de \timesfm{} 2.5:

\begin{lstlisting}[caption={Extracción de características con TimesFM (simplificado).}, label={lst:timesfm_extract}]
def extract_features(self, X, batch_size=32):
    N, seq_len, n_channels = X.shape
    point_preds = np.zeros((N, n_channels))
    quantile_preds = np.zeros((N, n_channels, 10))

    for start in range(0, N, batch_size):
        batch = X[start:end]
        flat_inputs = [batch[i, :, ch]
                       for i in range(B) for ch in range(n_channels)]

        point_pred, quant_pred = self._model.forecast(
            horizon=1, inputs=flat_inputs)
        # Des-aplanar resultados...
    return point_preds, quantile_preds
\end{lstlisting}

La diferencia clave con \chronos{} es que \timesfm{} genera directamente 10 cuantiles predefinidos en su salida nativa, sin necesidad de muestreo de Monte Carlo.


\subsection{Métricas probabilísticas: CRPS y cobertura}

Se implementan dos métricas probabilísticas estándar:

\paragraph{CRPS (Continuous Ranked Probability Score).}
Bajo el supuesto de distribución Normal, el CRPS se calcula analíticamente como \cite{gneiting2007scoring}:
\[
\text{CRPS}(\mu, \sigma, y) = \sigma \left[ z (2\Phi(z) - 1) + 2\phi(z) - \frac{1}{\sqrt{\pi}} \right], \quad z = \frac{y - \mu}{\sigma}
\]
donde $\Phi(\cdot)$ y $\phi(\cdot)$ son la función de distribución acumulada y la densidad de la Normal estándar, respectivamente.

El CRPS generaliza el MAE al espacio probabilístico: para predicciones deterministas ($\sigma \to 0$), el CRPS degenera al MAE estándar.

\paragraph{Cobertura del intervalo de predicción.}
La cobertura empírica mide la fracción de observaciones reales que caen dentro del intervalo de predicción:
\[
\text{Coverage}_{1-\alpha} = \frac{1}{N} \sum_{i=1}^{N} \mathbf{1}\left[ q_{\alpha/2}^{(i)} \leq y_i \leq q_{1-\alpha/2}^{(i)} \right]
\]
Un intervalo de predicción del 90\% bien calibrado debería tener una cobertura empírica cercana a 0.90.


% ======================================================================
\section{Módulo de Optimizadores (\code{fatigueset.models.optimizers})}
\label{sec:optimizers}
% ======================================================================

El módulo \code{optimizers.py} (477 líneas, 17.4 KB) extiende la búsqueda de hiperparámetros con Optuna al incluir la selección del \textbf{optimizador como hiperparámetro}.

\subsection{Arquitectura del módulo}

El módulo se estructura en cuatro componentes:

\begin{enumerate}
    \item \textbf{\code{build\_optimizer}}: Función de fábrica que instancia un optimizador de PyTorch a partir de su nombre y parámetros. Soporta Adam, AdamW, RMSprop y SGD (con momentum de Nesterov).

    \item \textbf{\code{sample\_optimizer\_params}}: Define el espacio de búsqueda condicional del optimizador en Optuna. Primero se selecciona el tipo de optimizador como variable categórica, y luego se activan únicamente los hiperparámetros relevantes:
    \begin{itemize}
        \item \textbf{Todos}: learning rate ($10^{-5}$ a $10^{-2}$, escala log) y weight decay ($10^{-7}$ a $10^{-3}$, escala log).
        \item \textbf{SGD solamente}: momentum (0.7 a 0.99).
        \item \textbf{RMSprop solamente}: alpha (0.9 a 0.999).
    \end{itemize}

    \item \textbf{Funciones \code{sample\_*} por familia}: Una función de muestreo por cada familia de modelos (LSTM, GRU, CNN-LSTM, TCN, Transformer, PatchTST, xLSTM, Random Forest) que define el espacio de búsqueda completo de la arquitectura \emph{incluyendo} el optimizador.

    \item \textbf{\code{sample\_model\_hyperparams}}: Función unificada que permite la búsqueda jerárquica y condicional entre familias de modelos en un único estudio de Optuna.
\end{enumerate}

\subsection{Búsqueda condicional con Optuna}

La búsqueda condicional es especialmente importante para evitar sugerir hiperparámetros irrelevantes (p.~ej., momentum para Adam), lo que reduce la dimensionalidad efectiva del espacio de búsqueda y mejora la eficiencia del sampler TPE.

\begin{lstlisting}[caption={Búsqueda condicional de optimizador en Optuna.}, label={lst:conditional_search}]
def sample_optimizer_params(trial, prefix):
    optimizer_name = trial.suggest_categorical(
        f"{prefix}_optimizer", ["Adam", "AdamW", "RMSprop", "SGD"])
    lr = trial.suggest_float(f"{prefix}_lr", 1e-5, 1e-2, log=True)
    weight_decay = trial.suggest_float(
        f"{prefix}_weight_decay", 1e-7, 1e-3, log=True)

    params = {"optimizer": optimizer_name, "lr": lr,
              "weight_decay": weight_decay}

    if optimizer_name == "SGD":
        params["momentum"] = trial.suggest_float(
            f"{prefix}_momentum", 0.7, 0.99)
    elif optimizer_name == "RMSprop":
        params["alpha"] = trial.suggest_float(
            f"{prefix}_alpha", 0.9, 0.999)
    return params
\end{lstlisting}

\subsection{Espacios de búsqueda por familia de modelo}

La Tabla~\ref{tab:search_spaces} resume los hiperparámetros buscados por Optuna para cada familia de modelo.

\begin{table}[H]
\centering
\caption{Resumen de hiperparámetros explorados por Optuna por familia de modelo.}
\label{tab:search_spaces}
\small
\begin{tabular}{@{}lp{9cm}@{}}
\toprule
\textbf{Familia} & \textbf{Hiperparámetros explorados} \\
\midrule
Random Forest & \code{n\_estimators} (50--300), \code{max\_depth}, \code{min\_samples\_split}, \code{min\_samples\_leaf}, \code{max\_features} \\
LSTM & \code{hidden\_size} (16--256), \code{num\_layers} (1--4), \code{dropout}, \code{batch\_size}, optimizador \\
GRU & \code{hidden\_size} (16--256), \code{num\_layers} (1--4), \code{dropout}, \code{batch\_size}, optimizador \\
CNN-LSTM & \code{conv\_channels}, \code{kernel\_size}, \code{pool\_size}, \code{hidden\_size}, \code{num\_layers}, \code{dropout}, optimizador \\
TCN & \code{num\_channels} (5 configuraciones), \code{kernel\_size} (2--8), \code{dropout}, optimizador \\
Transformer & \code{d\_model} (algebraicamente divisible), \code{nhead}, \code{num\_layers}, \code{dim\_feedforward}, \code{dropout}, optimizador \\
PatchTST & \code{patch\_len}, \code{stride}, \code{d\_model}, \code{n\_heads}, \code{num\_layers}, \code{dim\_ff}, \code{dropout}, optimizador \\
xLSTM & \code{hidden\_size} (32--256), \code{num\_layers} (1--4), \code{dropout}, optimizador \\
\bottomrule
\end{tabular}
\end{table}

\paragraph{Restricción de divisibilidad en Transformer y PatchTST.}
Para las arquitecturas basadas en auto-atención multi-cabeza, el parámetro $d_{\text{model}}$ debe ser divisible por el número de cabezas (\code{nhead}). En lugar de imponer esta restricción mediante pruning (descartando trials inválidos), se construye $d_{\text{model}} = \text{nhead} \times k$, donde $k$ es un multiplicador entero sugerido por Optuna. Esto garantiza la divisibilidad algebraicamente y evita desperdiciar presupuesto de búsqueda.


% ======================================================================
\section{Infraestructura de Computación en Servidor GPU}
\label{sec:infraestructura}
% ======================================================================

\subsection{Arquitectura del clúster UGR}

Los experimentos a escala completa se ejecutan en el clúster de computación GPU de la Universidad de Granada, con la siguiente configuración:

\begin{itemize}
    \item \textbf{Nodo cabecera}: \code{ngpu.ugr.es} --- punto de acceso SSH y gestor de colas SLURM.
    \item \textbf{Nodo de cómputo}: \code{titan} --- equipado con GPUs NVIDIA (capacidad reportada: 11.90 GiB por GPU).
    \item \textbf{Almacenamiento}: \code{/mnt/homeGPU/egullonl01/} --- almacenamiento persistente de 239 TB compartidos.
    \item \textbf{Gestor de recursos}: SLURM (\textit{Simple Linux Utility for Resource Management}).
\end{itemize}

\subsection{Entorno Conda remoto}

Se configuró un entorno Conda aislado en el almacenamiento GPU para evitar conflictos de dependencias:
\begin{verbatim}
/mnt/homeGPU/egullonl01/conda_tfg/
\end{verbatim}

Las dependencias principales instaladas incluyen: PyTorch (con CUDA), \code{momentfm}, \code{chronos-forecasting}, \code{timesfm}, \code{optuna}, \code{scikit-learn}, \code{pandas} y \code{matplotlib}.

\subsection{Scripts SLURM de ejecución}

Para posibilitar la ejecución desatendida y a escala de los experimentos, se estructuraron scripts bash controlados por SLURM. Estos scripts resuelven de forma nativa dos limitaciones críticas del entorno del clúster universitario:
\begin{enumerate}
    \item \textbf{Límite de Cuota de Disco en el Directorio Home}: Las librerías de HuggingFace y los servidores de Jupyter descargan por defecto los pesos de los modelos fundacionales (que suman más de 5~GB en total) y los archivos de sesión temporal en la carpeta de usuario de la cabecera (\code{\textasciitilde{}/.cache/}), la cual cuenta con una cuota restringida a 5~GB. Para evitar el fallo por desbordamiento (\emph{Disk quota exceeded}), se redireccionan los directorios de configuración y caché al almacenamiento de red persistente y con alta cuota del clúster mediante variables de entorno.
    \item \textbf{Liberación de Caché de GPU}: Dado que el recolector de PyTorch no devuelve de inmediato la memoria asignada al driver del sistema mientras el proceso de Python esté activo, se configuró la ejecución secuencial de notebooks en procesos independientes de Python, garantizando una GPU totalmente libre al inicio de cada modelo.
\end{enumerate}

El script \code{run\_foundation\_server.sh} ilustra esta configuración:

\begin{lstlisting}[language=bash, caption={Script SLURM completo para modelos fundacionales (\code{run\_foundation\_server.sh}).}, label={lst:slurm_script}]
#!/bin/bash
#SBATCH --job-name Foundation_TFG
#SBATCH --partition dios
#SBATCH --gres=gpu:1
#SBATCH --mem=32G
#SBATCH --time=12:00:00
#SBATCH --output=/mnt/homeGPU/egullonl01/tfg/Jupyters/slurm-%j.out
#SBATCH --error=/mnt/homeGPU/egullonl01/tfg/Jupyters/slurm-%j.out

export PATH="/opt/anaconda/anaconda3/bin:$PATH"
export PATH="/opt/anaconda/bin:$PATH"

# Redireccionar caches y configs para evitar problemas de cuota de disco
export IPYTHONDIR="/mnt/homeGPU/egullonl01/.ipython"
export JUPYTER_CONFIG_DIR="/mnt/homeGPU/egullonl01/.jupyter"
export JUPYTER_DATA_DIR="/mnt/homeGPU/egullonl01/.local/share/jupyter"
export JUPYTER_RUNTIME_DIR="/mnt/homeGPU/egullonl01/.local/share/jupyter/runtime"
export HF_HOME="/mnt/homeGPU/egullonl01/.cache/huggingface"
export XDG_CACHE_HOME="/mnt/homeGPU/egullonl01/.cache"

eval "$(conda shell.bash hook)"
conda activate /mnt/homeGPU/egullonl01/conda_tfg
cd /mnt/homeGPU/egullonl01/tfg/Jupyters

# Configurar SERVER_MODE de manera temporal en los notebooks
sed -i 's/SERVER_MODE = False/SERVER_MODE = True/g' 09_foundation_models_moment.ipynb
sed -i 's/SERVER_MODE = False/SERVER_MODE = True/g' 09_foundation_models_chronos.ipynb

# 1. Ejecutar MOMENT fine-tuning en un proceso aislado de Python
echo "Starting MOMENT..."
jupyter nbconvert --to notebook --execute \
    --ExecutePreprocessor.kernel_name=python3 \
    --inplace 09_foundation_models_moment.ipynb

# 2. Ejecutar Chronos zero-shot en su propio proceso limpio de Python
echo "Starting Chronos..."
jupyter nbconvert --to notebook --execute \
    --ExecutePreprocessor.kernel_name=python3 \
    --inplace 09_foundation_models_chronos.ipynb

# Restaurar estado original para Git
sed -i 's/SERVER_MODE = True/SERVER_MODE = False/g' 09_foundation_models_moment.ipynb
sed -i 's/SERVER_MODE = True/SERVER_MODE = False/g' 09_foundation_models_chronos.ipynb
\end{lstlisting}

\paragraph{Mecanismo \code{SERVER\_MODE}.}
Todos los notebooks incluyen una variable \code{SERVER\_MODE} que controla qué variante de modelo se carga. Los scripts SLURM la cambian automáticamente a \code{True} antes de la ejecución (para cargar los modelos grandes) y la revierten a \code{False} al finalizar (para mantener un estado limpio en Git). Los tres scripts comparten la configuración: partición \code{dios}, 1 GPU, 32 GB de RAM y 12 horas máximo. El script \code{run\_timesfm\_server.sh} sigue una estructura homóloga para procesar \timesfm{} de forma aislada.

\subsection{Pipeline de despliegue automático}

El script \code{scratch/deploy\_server.py} automatiza la sincronización de código entre el entorno local y el servidor remoto:

\begin{enumerate}
    \item Conexión SSH al nodo cabecera (\code{ngpu.ugr.es}).
    \item Subida de la librería completa \code{fatigueset-lib/} vía SFTP.
    \item Subida de los notebooks y scripts SLURM.
    \item Verificación de la integridad de los archivos transferidos.
\end{enumerate}

\subsection{Monitorización remota de jobs}

Se desarrollaron múltiples scripts de monitorización bajo \code{scratch/}:
\begin{itemize}
    \item \code{check\_slurm\_logs.py}: consulta \code{squeue} y muestra las últimas líneas de los logs de SLURM.
    \item \code{check\_slurm\_job\_details.py}: ejecuta \code{scontrol show job} para obtener el estado detallado.
    \item \code{inspect\_foundation\_notebook.py}: descarga el notebook ejecutado del servidor y extrae las salidas de cada celda para diagnóstico remoto.
    \item \code{inspect\_remote\_output.py}: lista los archivos generados en el directorio \code{output/} del servidor.
\end{itemize}


% ======================================================================
\section{Resultados Experimentales}
\label{sec:resultados}
% ======================================================================


\subsection{Tabla comparativa global}

La Tabla~\ref{tab:resultados_globales} presenta los resultados de todos los modelos evaluados. Las métricas son MAE (error absoluto medio), RMSE (raíz del error cuadrático medio) y $R^2$ (coeficiente de determinación). Para los modelos fundacionales con predicciones separadas por dimensión, se reportan las dos dimensiones explícitamente.

\paragraph{Nota metodológica sobre los modelos clásicos.}
Es sumamente importante destacar una discrepancia metodológica en los resultados de los modelos clásicos de Machine Learning. Mientras que las métricas de los modelos de Deep Learning (personalizados y fundacionales) corresponden a la evaluación sobre el conjunto de validación en validación cruzada (reflejando capacidad de generalización real), el MAE y RMSE de los modelos clásicos se evaluaron sobre el propio conjunto de entrenamiento. Esto introduce un sesgo optimista severo y explica sus errores aparentemente reducidos (p.~ej., un MAE de 0.88 para Random Forest o 0.00 para el árbol de decisión). El $R^2$ negativo en la validación cruzada para estos mismos modelos clásicos confirma que sufren de un sobreajuste severo y carecen de capacidad real de generalización sobre sujetos no vistos.

\begin{table}[H]
\centering
\caption{Resultados comparativos globales y tiempos de cómputo de todos los modelos del TFG.}
\label{tab:resultados_globales}
\small
\begin{tabular}{@{}lcccccc@{}}
\toprule
\textbf{Modelo} & \textbf{Tipo} & \textbf{MAE Medio} & \textbf{RMSE Medio} & \textbf{$R^2$ Medio} & \textbf{Params} & \textbf{Tiempo Total} \\
\midrule
\multicolumn{7}{l}{\textit{Modelos clásicos de ML (validación cruzada KFold para $R^2$, entrenamiento para MAE/RMSE)$^\S$}} \\
Random Forest & Ensamble & 0.88$^\S$ & 1.04$^\S$ & $-$2.41 & --- & $<$15 s \\
KNN (k=3) & Instancias & 1.69$^\S$ & 2.20$^\S$ & $-$3.57 & --- & $<$5 s \\
Decision Tree & Árbol & 0.00$^*$ & 0.00$^*$ & 1.00$^*$ & --- & $<$2 s \\
\midrule
\multicolumn{7}{l}{\textit{Modelos de Deep Learning custom (MAE/RMSE/R² promediados sobre ambas dimensiones, Optuna v2 CV)}} \\
\textbf{Custom GRU (Opt. v2)} & Recurrente & \textbf{16.99} & \textbf{20.14} & $-$0.62 & 568k / 58k & 7.057 s ($\approx$ 1,96 h) \\
\textbf{Custom xLSTM (Opt. v2)} & Recurrente & \textbf{17.33} & \textbf{20.50} & $-$0.69 & 625k / 287k & 12.866 s ($\approx$ 3,57 h) \\
\textbf{Custom LSTM (Opt. v2)} & Recurrente & \textbf{17.50} & \textbf{20.45} & $-$0.68 & 887k / 120k & 3.045 s ($\approx$ 50,7 min) \\
Custom CNN-LSTM (Opt. v2) & Híbrido & 17.56 & 20.62 & $-$0.71 & 84k / 807k & 3.183 s ($\approx$ 53,0 min) \\
Custom TCN (Opt. v2) & Convolucional & 17.70 & 21.52 & $-$0.96 & 175k / 134k & 298 s ($\approx$ 5,0 min) \\
Custom Transformer (Opt. v2) & Transformer & 18.25 & 21.10 & $-$0.80 & 396k / 130k & 767 s ($\approx$ 12,8 min) \\
Custom PatchTST (Opt. v2) & Transformer & 19.44 & 22.38 & $-$1.04 & 77k / 655k & 862 s ($\approx$ 14,4 min) \\
Custom RNN (Baseline) & Recurrente & 30.09 & 34.93 & $-$5.07 & 1.890 & 18,5 s \\
\midrule
\multicolumn{7}{l}{\textit{Modelos fundacionales (métricas de validación por dimensión, zero-shot/fine-tuning)}} \\
\multirow{2}{*}{\textbf{Chronos (chronos-t5-base)}} & \multirow{2}{*}{Foundation} & Fís: 14.11 & Fís: 17.03 & Fís: $-$0.45 & \multirow{2}{*}{710M$^\dagger$} & \multirow{2}{*}{1.240 s ($\approx$ 20,7 min)} \\
& & Men: 18.95 & Men: 22.22 & Men: $-$1.17 & & \\
\multirow{2}{*}{\textbf{\timesfm{} 2.5}} & \multirow{2}{*}{Foundation} & Fís: 16.39 & Fís: 20.15 & Fís: $-$1.10 & \multirow{2}{*}{200M$^\dagger$} & \multirow{2}{*}{536,7 s ($\approx$ 8,9 min)} \\
& & Men: 22.38 & Men: 26.71 & Men: $-$1.43 & & \\
\multirow{2}{*}{\moment{} (large)} & \multirow{2}{*}{Foundation} & Fís: 24.64 & Fís: 28.98 & Fís: $-$3.88 & \multirow{2}{*}{2.050$^\ddagger$} & \multirow{2}{*}{5.210 s ($\approx$ 86,8 min)} \\
& & Men: 36.68 & Men: 41.37 & Men: $-$8.07 & & \\
\bottomrule
\end{tabular}
\begin{flushleft}
\footnotesize
$^*$ Decision Tree con MAE=0 e $R^2=1.00$ indica sobreajuste severo (memorización completa del conjunto de entrenamiento). \\
$^\dagger$ Parámetros totales del backbone preentrenado (0 parámetros entrenados, linear probe externo). \\
$^\ddagger$ Parámetros entrenables de la cabeza de regresión (backbone de 341M congelado). \\
$^\S$ Métricas de MAE y RMSE calculadas sobre el conjunto de entrenamiento (no de validación cruzada).
\end{flushleft}
\end{table}


\subsection{MOMENT: resultados por fold}

La Tabla~\ref{tab:moment_folds} muestra los resultados de \moment{}-1-large con fine-tuning de la cabeza de regresión, usando 5-fold GroupKFold por participante.

\begin{table}[H]
\centering
\caption{Resultados de MOMENT-1-large por fold (GroupKFold, 5 splits, 30 épocas).}
\label{tab:moment_folds}
\small
\begin{tabular}{@{}ccccccc@{}}
\toprule
\textbf{Fold} & \textbf{MAE Fís.} & \textbf{RMSE Fís.} & \textbf{$R^2$ Fís.} & \textbf{MAE Men.} & \textbf{RMSE Men.} & \textbf{$R^2$ Men.} \\
\midrule
1 & 26.30 & 30.19 & $-$3.08 & 29.38 & 31.41 & $-$7.00 \\
2 & 24.82 & 29.11 & $-$2.66 & 44.46 & 48.61 & $-$5.11 \\
3 & 33.21 & 36.70 & $-$4.52 & 52.77 & 56.53 & $-$6.78 \\
4 & 25.49 & 26.95 & $-$8.55 & 27.42 & 28.08 & $-$20.56 \\
5 & 13.37 & 21.93 & $-$0.59 & 29.37 & 42.23 & $-$0.90 \\
\midrule
\textbf{Media} & \textbf{24.64} & \textbf{28.98} & \textbf{$-$3.88} & \textbf{36.68} & \textbf{41.37} & \textbf{$-$8.07} \\
$\pm$ Std & $\pm$6.39 & $\pm$4.80 & $\pm$2.65 & $\pm$10.12 & $\pm$10.57 & $\pm$6.62 \\
\bottomrule
\end{tabular}
\begin{flushleft}
\footnotesize
Configuración: batch size 32, learning rate $3 \times 10^{-5}$, early stopping con paciencia 10. \\
Tiempo total de fine-tuning: 5.209,7 s ($\approx$ 87 min).
\end{flushleft}
\end{table}

\paragraph{Observaciones.}
\begin{itemize}
    \item El Fold 5 muestra el mejor rendimiento ($R^2_\text{fís} = -0.59$, $R^2_\text{men} = -0.90$), sugiriendo que los participantes de validación de ese fold tienen patrones más alineados con los del entrenamiento.
    \item El Fold 4 presenta el peor $R^2$ mental ($-20.56$), indicando una distribución de fatiga mental muy diferente en ese grupo de participantes.
    \item La alta varianza inter-fold ($\pm$6.62 en $R^2$ mental) refleja la heterogeneidad individual inherente a \fatigueset{}.
\end{itemize}


\subsection{Chronos: métricas zero-shot y comparación de escalas}

La Tabla~\ref{tab:chronos_metricas} presenta los resultados probabilísticos y deterministas de \chronos{} utilizando las variantes \code{chronos-t5-tiny} (8M, ejecutado en local) y \code{chronos-t5-base} (710M, ejecutado en servidor).

\begin{table}[H]
\centering
\caption{Métricas probabilísticas y deterministas de las variantes de Chronos.}
\label{tab:chronos_metricas}
\small
\begin{tabular}{@{}lcccc@{}}
\toprule
\multirow{2}{*}{\textbf{Métrica}} & \multicolumn{2}{c}{\textbf{Chronos-T5-tiny (8M)}} & \multicolumn{2}{c}{\textbf{Chronos-T5-base (710M)}} \\
\cmidrule(lr){2-3} \cmidrule(lr){4-5}
& \textbf{Fatiga Física} & \textbf{Fatiga Mental} & \textbf{Fatiga Física} & \textbf{Fatiga Mental} \\
\midrule
MAE & 17.25 & 23.40 & 14.11 & 18.95 \\
RMSE & 21.05 & 27.60 & 17.03 & 22.22 \\
CRPS (media $\pm$ std) & 12.59 $\pm$ 3.98 & 17.28 $\pm$ 6.30 & 10.17 $\pm$ 2.63 & 13.69 $\pm$ 4.88 \\
Cobertura 90\% & 68.4\% & 63.2\% & 78.9\% & 73.8\% \\
\bottomrule
\end{tabular}
\end{table}


\subsection{TimesFM: métricas puntuales y probabilísticas}

\begin{table}[H]
\centering
\caption{Resultados de TimesFM 2.5 (200M) con zero-shot + linear probe.}
\label{tab:timesfm_resultados}
\small
\begin{tabular}{@{}lcc@{}}
\toprule
\textbf{Métrica} & \textbf{Fatiga Física} & \textbf{Fatiga Mental} \\
\midrule
MAE & 16.39 & 22.38 \\
RMSE & 20.15 & 26.71 \\
$R^2$ & $-$1.10 & $-$1.43 \\
\midrule
CRPS (media $\pm$ std) & 12.53 $\pm$ 3.93 & 17.30 $\pm$ 6.35 \\
Cobertura 90\% & 68.7\% & 62.8\% \\
\midrule
Parámetros entrenados & 0 (probe externo) & \\
Tiempo total & 536.7 s & \\
\bottomrule
\end{tabular}
\end{table}

\paragraph{Observaciones.}
\begin{itemize}
    \item \timesfm{} alcanza el mejor MAE global del proyecto (16.39 en fatiga física, 22.38 en mental), superando a los modelos de DL custom y a MOMENT.
    \item Las métricas probabilísticas de \timesfm{} son muy similares a las de \chronos{}-tiny (CRPS física: 12.53 vs 12.59; cobertura 90\% física: 68.7\% vs 68.4\%), sugiriendo que ambos modelos producen distribuciones predictivas de calidad comparable.
    \item La cobertura del 90\% está por debajo del nominal (68--69\% en lugar del esperado 90\%), lo que indica que los intervalos de predicción son demasiado estrechos (subconfianza).
\end{itemize}


\subsection{Optuna: resultados del experimento extendido con optimizadores personalizados (Optuna v2)}
\label{sec:optuna_v2_resultados}

El experimento de búsqueda bayesiana extendida (\textbf{Optuna v2}) integró la selección del optimizador como hiperparámetro categórico (\code{Adam}, \code{AdamW}, \code{RMSprop}, \code{SGD}) junto a las tasas de aprendizaje, aceleración de momento, \textit{weight decay} e hiperparámetros estructurales específicos de cada arquitectura ($h$, $L$, $d_{\text{model}}$, $p$, etc.). La búsqueda se ejecutó en el clúster GPU con un presupuesto de $N_{\text{trials}} = 25$, validación cruzada 3-fold GroupKFold por participante y 10 épocas por trial.

La Tabla~\ref{tab:optuna_v2_resumen} detalla los mejores hiperparámetros seleccionados por Optuna y los errores resultantes para cada objetivo y familia de modelos.

\begin{table}[H]
\centering
\caption{Mejores hiperparámetros, métricas y tiempos de cómputo alcanzados por Optuna v2 por variable objetivo.}
\label{tab:optuna_v2_resumen}
\small
\begin{tabular}{@{}llccccc@{}}
\toprule
\textbf{Objetivo} & \textbf{Modelo} & \textbf{Optimizador} & \textbf{Learning Rate} & \textbf{MAE Medio} & \textbf{RMSE Medio} & \textbf{Tiempo Búsqueda} \\
\midrule
\multirow{7}{*}{\textbf{Fatiga Física}} 
& Custom GRU & Adam & $1.00 \times 10^{-3}$ & \textbf{13.94} & \textbf{16.86} & $<$1 s \\
& Custom xLSTM & RMSprop & $6.34 \times 10^{-4}$ & 14.46 & 17.54 & 1.870 s (31,2 min) \\
& Custom TCN & RMSprop & $1.12 \times 10^{-4}$ & 14.79 & 18.27 & 6,0 s \\
& Custom LSTM & SGD & $1.44 \times 10^{-3}$ & 15.06 & 17.71 & $<$1 s \\
& Custom CNN-LSTM & AdamW & $1.09 \times 10^{-3}$ & 15.10 & 17.91 & 90,4 s (1,5 min) \\
& Custom PatchTST & AdamW & $8.90 \times 10^{-5}$ & 15.68 & 17.87 & $<$1 s \\
& Custom Transformer & SGD & $4.63 \times 10^{-4}$ & 15.79 & 18.31 & 17,0 s \\
\midrule
\multirow{7}{*}{\textbf{Fatiga Mental}} 
& Custom LSTM & Adam & $1.70 \times 10^{-3}$ & \textbf{19.93} & \textbf{23.19} & 3.045 s (50,7 min) \\
& Custom CNN-LSTM & AdamW & $8.96 \times 10^{-4}$ & 20.02 & 23.32 & 3.093 s (51,5 min) \\
& Custom GRU & Adam & $1.59 \times 10^{-3}$ & 20.04 & 23.41 & 7.057 s (117,6 min) \\
& Custom xLSTM & SGD & $2.06 \times 10^{-3}$ & 20.19 & 23.45 & 10.995 s (183,3 min) \\
& Custom TCN & Adam & $3.10 \times 10^{-4}$ & 20.60 & 24.77 & 292 s (4,9 min) \\
& Custom Transformer & RMSprop & $2.11 \times 10^{-3}$ & 20.70 & 23.89 & 750 s (12,5 min) \\
& Custom PatchTST & SGD & $1.98 \times 10^{-5}$ & 23.19 & 26.88 & 862 s (14,4 min) \\
\bottomrule
\end{tabular}
\end{table}

\paragraph{Observaciones del experimento de optimización.}
\begin{itemize}
    \item \textbf{Disparidad de Optimizadores}: Para arquitecturas recurrentes puras (\code{GRU} y \code{LSTM}), optimizadores basados en momento adaptativo (\code{Adam}) alcanzaron la convergencia más rápida. Para \code{xLSTM} y \code{TCN} en fatiga física, \code{RMSprop} ofreció una estabilidad superior en las puertas exponenciales y convoluciones dilatadas.
    \item \textbf{Diferencias entre objetivos}: La fatiga física es sustancialmente más predecible ($\text{MAE} \approx 13.94$ -- $15.79$) que la fatiga mental ($\text{MAE} \approx 19.93$ -- $23.19$), lo que se explica por la mayor correlación directa de las señales cardiorrespiratorias y de acelerometría con el esfuerzo neuromuscular.
\end{itemize}

\subsection{Análisis de Eficiencia Computacional y Tiempos de Cómputo}
\label{sec:tiempos_computo}

Un aspecto fundamental en la evaluación empírica de este Trabajo de Fin de Grado es el intercambio (\textit{trade-off}) entre la precisión predictiva alcanzada y el coste computacional requerido por cada paradigma:

\begin{enumerate}
    \item \textbf{Machine Learning Clásico ($<15$ segundos)}: Los modelos sobre datos tabulares agregados (Random Forest, KNN, Árboles) presentan un coste de entrenamiento prácticamente despreciable. Sin embargo, sufren de un sobreajuste severo debido a la pérdida del contexto secuencial continuo.
    
    \item \textbf{Modelos Fundacionales Zero-Shot (8,9 a 20,7 minutos)}:
    \begin{itemize}
        \item \textbf{TimesFM 2.5 (200M)} completó su evaluación global en solo \textbf{536,7 s ($\approx 8,9$ min)}, ofreciendo una de las mejores relaciones entre rapidez e inferencia determinista/probabilística.
        \item \textbf{Chronos-t5-base (710M)} requirió \textbf{1.240 s ($\approx 20,7$ min)} para muestrear 100 trayectorias autorregresivas por ventana. A cambio de este modesto tiempo de cómputo sin entrenamiento de pesos, alcanzó las mejores métricas absolutas del proyecto ($\text{MAE} = 14,11$).
    \end{itemize}
    
    \item \textbf{Fine-Tuning de Modelo Fundacional ($\approx 86,8$ minutos)}:
    El ajuste de la cabeza lineal de \textbf{MOMENT-1-large} requirió \textbf{5.210 s ($\approx 86,8$ min)} a lo largo de 5 splits de validación cruzada. El elevado tiempo proviene del paso forward de 341M de parámetros para extraer representaciones de parches, sin lograr superar a los modelos zero-shot o redes custom optimizadas.
    
    \item \textbf{Modelos Deep Learning Custom con Optuna v2 (5 min a 3,57 horas por familia)}:
    La búsqueda bayesiana jerárquica con validación cruzada acumuló un tiempo total de computación de aproximadamente **8,8 horas de GPU Titan** en el clúster universitario:
    \begin{itemize}
        \item Las arquitecturas ligeras convolucionales y de atención temporal (\code{TCN}, \code{Transformer}, \code{PatchTST}) resultaron muy eficientes, completando sus búsquedas bayesianas completas de 50 trials en **5 a 14 minutos**.
        \item Las arquitecturas recurrentes de alta capacidad como \code{xLSTM} (10.995 s $\approx 3,05$ h para fatiga mental) y \code{GRU} (7.057 s $\approx 1,96$ h) requirieron un esfuerzo computacional significativamente mayor debido a los bucles secuenciales unrolled paso a paso en memoria GPU.
    \end{itemize}
\end{enumerate}


\subsection{Discusión y análisis comparativo}

Los resultados experimentales revelan varias conclusiones clave:

\begin{enumerate}
    \item \textbf{\chronos{} (chronos-t5-base) es el mejor modelo del proyecto}, logrando el menor error absoluto medio tanto para fatiga física (MAE de 14.11) como para fatiga mental (MAE de 18.95). Su paradigma zero-shot, que discretiza los valores numéricos mediante cuantiles y utiliza la potente base lingüística del T5, demuestra ser una excelente herramienta para extraer representaciones fisiológicas universales y robustas.
    
    \item \textbf{\timesfm{} 2.5 (200M) muestra también un rendimiento muy competitivo} (MAE físico: 16.39; mental: 22.38), superando a todos los modelos recurrentes. Esto confirma que los decoders de series temporales basados en parches o tokens discretos son superiores a los encoders o arquitecturas secuenciales clásicas.
    
    \item \textbf{Los modelos de DL custom optimizados con Optuna} (LSTM, GRU, PatchTST) logran mejorar significativamente frente a sus baselines no optimizados. El MAE promedio de LSTM y GRU se redujo de ~30 a ~16.8 - 17.0 tras la optimización bayesiana de la tasa de aprendizaje y dropout, igualando la precisión del modelo TCN (MAE ~17.6) y PatchTST (MAE ~17.5).
    
    \item \textbf{El fine-tuning de cabeza} (\moment{}) no mejora sobre el zero-shot en datasets pequeños. Con solo 660 ventanas y alta variabilidad individual, el fine-tuning de MOMENT (large) es vulnerable al sobreajuste (MAE físico: 24.64, mental: 36.68), rindiendo peor que los modelos zero-shot o redes custom optimizadas.
    
    \item \textbf{Los valores de $R^2$ negativos} en casi todos los modelos indican que la predicción individual de fatiga sigue siendo un problema sumamente complejo debido a la gran variabilidad inter-sujeto de \fatigueset{}. No obstante, la reducción sustancial del MAE por parte de Chronos (14.11 frente a la desviación estándar de la fatiga física de 16.52) representa un avance experimental significativo.
    
    \item \textbf{La calibración probabilística de Chronos-base es notablemente superior}: logra coberturas del intervalo de predicción del 78.9\% (física) y 73.8\% (mental), un incremento considerable frente al 68.7\%/62.8\% de TimesFM o el 68.4\%/63.2\% de Chronos-tiny, evidenciando el efecto positivo de la escala del modelo en la estimación de la incertidumbre.
\end{enumerate}


% ======================================================================
\section{Problemas Técnicos y Soluciones}
\label{sec:problemas}
% ======================================================================

\subsection{CUDA Out of Memory: diagnóstico y corrección}
\label{sec:problemas_oom}

\paragraph{Problema.}
Al ejecutar el notebook \code{09\_foundation\_models.ipynb} en el servidor, \moment{}-1-large se entrenó correctamente en los 5 folds de validación cruzada, pero \chronos{}-large falló con el error:

\begin{verbatim}
CUDA out of memory. Tried to allocate 2.94 GiB. GPU 0 has a
total capacity of 11.90 GiB of which 210.81 MiB is free.
Of the allocated memory 11.50 GiB is allocated by PyTorch.
\end{verbatim}

\paragraph{Causa raíz.}
La GPU asignada por SLURM tiene 11.90 GiB de capacidad (no 24 GiB como se esperaba inicialmente). Tras el fine-tuning de \moment{}, PyTorch retenía 11.50 GiB de tensores en la GPU. Aunque el backbone compartido se eliminó explícitamente al final de la función (\code{del backbone\_shared; gc.collect(); torch.cuda.empty\_cache()}), las activaciones intermedias del forward pass permanecían en memoria.

\paragraph{Análisis técnico.}
El problema subyacente es que en el bucle de entrenamiento:
\begin{lstlisting}
model.train()
for xb, yb in train_loader:
    pred = model(xb)      # Forward pass sin torch.no_grad()
    loss = loss_fn(pred, yb)
    loss.backward()
\end{lstlisting}
Aunque el backbone tiene \code{requires\_grad=False}, al ejecutarse en modo \code{model.train()} sin \code{torch.no\_grad()} alrededor de la llamada al backbone, PyTorch almacena las activaciones intermedias de las 24 capas Transformer para el grafo de autodiferenciación. Estas activaciones consumen varios GiB de memoria GPU.

\paragraph{Solución implementada.}
Se implementó una estrategia de \textbf{backbone compartido} para reducir el número de cargas del modelo de 5 (una por fold) a 1:
\begin{lstlisting}[caption={Backbone compartido en finetune\_moment\_kfold.}]
# Cargar backbone UNA sola vez
backbone_shared = MOMENTPipeline.from_pretrained(checkpoint, ...)

for fold_idx, (train_idx, val_idx) in enumerate(kf.split(...)):
    model = MOMENTFatigueRegressor(
        ..., backbone=backbone_shared)  # Reutilizar
    model = model.to(device)
    # ... entrenamiento ...
    del model, optimizer
    gc.collect(); torch.cuda.empty_cache()

# Liberar backbone al final
del backbone_shared
gc.collect(); torch.cuda.empty_cache()
\end{lstlisting}

\paragraph{Soluciones finales aplicadas.}
Se han implementado tres soluciones complementarias para resolver los problemas de memoria GPU:
\begin{enumerate}
    \item \textbf{Inferencia sin gradientes}: Se envolvió el paso forward del backbone congelado en un bloque \code{with torch.no\_grad():} en \code{MOMENTFatigueRegressor.forward} (cuando \code{freeze\_backbone} es \code{True}). Esto evitó de forma efectiva que PyTorch retuviera los grafos de activaciones en memoria.
    \item \textbf{Separación de procesos secuenciales}: Dado que el distribuidor de caché de PyTorch (\textit{caching allocator}) mantiene bloques asignados para acelerar futuras peticiones y no los devuelve al driver de la GPU mientras el proceso de Python esté vivo, se dividió el notebook \code{09\_foundation\_models.ipynb} en dos notebooks independientes: \code{09\_foundation\_models\_moment.ipynb} y \code{09\_foundation\_models\_chronos.ipynb}. El script de ejecución de SLURM ejecuta ambos de forma secuencial en procesos de Python independientes. Al terminar el primer notebook, su proceso se destruye liberando por completo la memoria de la GPU, lo que permite al segundo notebook ejecutarse sobre una GPU totalmente vacía.
    \item \textbf{Aislamiento de TimesFM}: Al igual que Chronos y MOMENT, el modelo de Google \timesfm{} 2.5 (200M) se ejecuta en su propio notebook independiente \code{10\_timesfm.ipynb} y mediante el script de lanzamiento \code{run\_timesfm\_server.sh}. Esto garantiza un entorno limpio y predecible de memoria, libre de cualquier tipo de fragmentación previa del caching allocator.
\end{enumerate}

\subsection{Exceso de Cuota de Disco en Directorio Home (Disk Quota Exceeded)}
\label{sec:problemas_cuota}

\paragraph{Problema.}
Durante la primera descarga de los backbones fundacionales preentrenados, la ejecución abortaba con un error de escritura en disco:
\begin{verbatim}
OSError: [Errno 122] Disk quota exceeded
\end{verbatim}

\paragraph{Causa.}
Por defecto, la librería \code{transformers} de HuggingFace descarga y almacena los pesos de los modelos (MOMENT-large: $\sim$1.5~GB, Chronos-base: $\sim$2.8~GB, TimesFM: $\sim$1~GB) en el directorio personal del usuario bajo el host cabecera (\code{\textasciitilde{}/.cache/huggingface}). La cuota de almacenamiento para directorios de usuario del servidor universitario está limitada a 5~GB, lo que bloqueaba la descarga concurrente o consecutiva de los modelos fundacionales.

\paragraph{Solución implementada.}
Se redirigieron de forma explícita todas las variables de entorno de caché, ejecución y configuración al almacenamiento compartido persistente del clúster de computación, el cual posee una capacidad de almacenamiento de 239~TB y sin restricciones estrictas de cuota individual. En el inicio de los scripts SLURM, se establecen las siguientes redirecciones:
\begin{itemize}
    \item \code{HF\_HOME}: Redirigido a \code{/mnt/homeGPU/egullonl01/.cache/huggingface} para almacenar de forma segura los pesos de HuggingFace.
    \item \code{XDG\_CACHE\_HOME}: Redirigido a \code{/mnt/homeGPU/egullonl01/.cache} para el almacenamiento de caché general.
    \item \code{IPYTHONDIR}: Redirigido a \code{/mnt/homeGPU/egullonl01/.ipython} para la caché de IPython.
    \item \code{JUPYTER\_*}: Las variables de configuración y datos de Jupyter (\code{JUPYTER\_CONFIG\_DIR}, \code{JUPYTER\_DATA\_DIR}, \code{JUPYTER\_RUNTIME\_DIR}) se redirigieron a subcarpetas equivalentes en \code{/mnt/homeGPU/egullonl01/}.
\end{itemize}
Esta reconfiguración evitó por completo los desbordamientos de cuota en el home de usuario, garantizando descargas de modelos sin fricciones.

\subsection{Exceso de tiempo en SLURM (Optuna)}
\label{sec:problemas_slurm}

\paragraph{Problema.}
El Job 155622 (\code{experimento\_optuna\_optimizadores.ipynb}) fue cancelado por SLURM tras exceder las 12 horas configuradas:
\begin{verbatim}
slurmstepd: error: *** JOB 155622 ON titan CANCELLED AT
2026-07-03T19:52:18 DUE TO TIME LIMIT ***
\end{verbatim}

\paragraph{Causa.}
El espacio de búsqueda del experimento con optimizadores personalizados es significativamente más grande que el experimento básico (8 familias $\times$ 4 optimizadores $\times$ múltiples hiperparámetros condicionales). Cada trial requiere entrenar un modelo completo con validación cruzada GroupKFold (3 splits), lo que multiplica el tiempo de ejecución por 3.

\paragraph{Solución aplicada.}
Para mitigar este problema sin comprometer el rigor experimental, se optimizó el presupuesto de búsqueda reduciendo los hiperparámetros de escala del servidor a $N_{\text{trials}} = 25$, $N_{\text{folds}} = 3$ y $\text{épocas}_{\text{trial}} = 10$. Este ajuste reduce el total de épocas acumuladas de entrenamiento a unas 6.300, disminuyendo la estimación de tiempo de ejecución a aproximadamente 6,5 horas de GPU. El experimento fue relanzado con éxito bajo esta nueva configuración, completándose holgadamente dentro del límite de 12 horas.


% ======================================================================
\section{Conclusiones y Trabajo Futuro}
\label{sec:conclusiones}
% ======================================================================

\subsection{Comparativa final de paradigmas}

\begin{enumerate}
    \item \textbf{Modelos fundacionales zero-shot} (\timesfm{}, \chronos{}) alcanzan resultados competitivos sin necesidad de entrenamiento específico, validando su utilidad como extractores de características universales para dominios no vistos.

    \item \textbf{El fine-tuning de cabeza} (\moment{}) no mejora necesariamente sobre el zero-shot en datasets pequeños. Con solo 660 ventanas y alta variabilidad individual, el fine-tuning es vulnerable al sobreajuste a pesar del pequeño número de parámetros entrenables (2.050).

    \item \textbf{Los modelos custom bien diseñados} (TCN, PatchTST) siguen siendo competitivos, especialmente cuando la arquitectura se adapta específicamente al problema (convoluciones dilatadas para señales fisiológicas, parches para capturas temporales multi-escala).

    \item \textbf{La heterogeneidad individual} de \fatigueset{} representa el mayor desafío experimental, como evidencia la alta varianza inter-fold y los valores negativos de $R^2$ en todos los modelos.
\end{enumerate}

\subsection{Tareas pendientes}

\begin{enumerate}
    \item \textbf{Optuna con optimizadores (Completado)}: Se redujo el presupuesto de búsqueda a $N_{\text{trials}} = 25$ y $\text{épocas} = 10$, y se relanzó con éxito el experimento extendido dentro del límite de 12 horas, persistiendo los hiperparámetros óptimos resultantes.
    \item \textbf{Calibración probabilística avanzada}: explorar técnicas de calibración post-hoc (Platt scaling, regresión isotónica) para elevar la cobertura empírica al 90\% nominal.
    \item \textbf{Análisis por participante}: evaluar métricas condicionadas al individuo para cuantificar la variabilidad inter-sujeto y proponer modelos personalizados de regresión.
\end{enumerate}


% ======================================================================
% BIBLIOGRAFÍA
% ======================================================================

\begin{thebibliography}{99}

\bibitem{kalanadhabhatta2022fatigueset}
Kalanadhabhatta, M., Zou, A., Rajagopalan, A., et al. (2022).
\textit{FatigueSet: A Multi-modal Dataset for Modeling Mental and Physical Fatigue}.
arXiv:2203.06924.

\bibitem{kottapalli2025survey}
Kottapalli, S. R. K., Hubli, K., Chandrashekhara, S., et al. (2025).
\textit{Foundation Models for Time Series: A Survey}.
arXiv.

\bibitem{beck2024xlstm}
Beck, M., Pöppel, K., Spanring, M., et al. (2024).
\textit{xLSTM: Extended Long Short-Term Memory}.
arXiv:2405.20455.

\bibitem{goswami2024moment}
Goswami, M., Szafer, K., Choudhry, A., et al. (2024).
\textit{MOMENT: A Family of Open Time-series Foundation Models}.
International Conference on Machine Learning (ICML).

\bibitem{ansari2024chronos}
Ansari, A. F., Stella, L., Turkmen, C., et al. (2024).
\textit{Chronos: Learning the Language of Time Series}.
arXiv:2403.07815.

\bibitem{das2024timesfm}
Das, A., Kong, W., Leach, A., et al. (2024).
\textit{A Decoder-only Foundation Model for Time-series Forecasting}.
International Conference on Machine Learning (ICML).

\bibitem{gneiting2007scoring}
Gneiting, T. \& Raftery, A. E. (2007).
\textit{Strictly Proper Scoring Rules, Prediction, and Estimation}.
Journal of the American Statistical Association, 102(477), 359--378.

\end{thebibliography}

\end{document}
